% =============================================================================
%  DRAFT for 05-results.tex (not yet pasted; the Overleaf review round is open)
% =============================================================================
% Proposed as a new subsection after \cref{sec:res-announced}: an EXTENSION and robustness check
% on the 200 bp cap, which stays the paper's policy. The cap has not been solved as a known rule.
% PROVISIONAL NUMBERS (2026-11-13), from the global solution of PR #38 (Bohn coefficient 0.25;
% deposits clear nationally in the steady state but across the union in the dynamics), with
% the proportional rule at k = 2 solved as an announced rule (branch new-rules). Collocation
% residuals at the grid points are 1.2e-8. No claim of global convergence: the impact quarter
% is cleared exactly, with continuation values from the fitted decision rules, and the last
% paragraph reports how the answers change when two quarters are cleared jointly.
% Accounting decompositions (exact identities) and controlled counterfactuals (experiments
% that are not policy equilibria) are labelled as such in the text.
% Sources: Claude files/diagnostics/bft_rule/FINDINGS.md, sections "Forensic decomposition of
% k = 2" and "Why the price of capital moves (part 3)";
% results/GLOBAL_2026-11-12_pr38_bft/forensic/*.json.

\subsection{An announced proportional purchase rule}
\label{sec:res-proportional}

A spread cap that agents know about is hard to compute, because near the top of the risk box
the cap must hold against a fair spread of more than 500\,bp. As a tractable announced
counterpart we study the proportional rule of \citet{bi2026asset}. The central bank
buys $D$ bonds whenever their expected excess return exceeds its value at the rest point,
\begin{equation}
q^{D}_{t}\,M_{t+1}=k\,\bar q^{D}\bar B_{D}\,\bigl(\mathit{ER}_{t}-\overline{\mathit{ER}}\bigr),
\qquad M_{t+1}\ge M_{t},
\label{eq:proportional-rule}
\end{equation}
where $M_{t+1}$ is its holding of $D$ bonds, $\mathit{ER}_t$ the quarterly expected excess
return of the $D$ bond over the $D$ deposit rate, $\overline{\mathit{ER}}=15$\,bp a year its
value at the rest point, and $\bar q^{D}\bar B_D$ the market value of the $D$ debt at the
steady state. The rule never sells and is known to all agents. We set the response coefficient
$k=2$, against 16 in \citet{bi2026asset}.

At the rest point the promise does almost nothing: the sovereign spread falls by 2\,bp and
output rises by 0.04\%. At the headline shock it changes the outcome without moving the bond
market much. \Cref{tab:proportional} shows that the rule buys 1.1\% of the $D$ debt and lowers
the sovereign spread by only 9\,bp, yet the credit spread $r^{wc}_{D,t}-r_{D,t}$ falls from
70\,bp to zero: the incentive constraint goes slack. At a 6\% default probability the rule buys
4.7\% of the debt, the sovereign spread falls by 43\,bp and the credit spread by 305\,bp. The
credit spread responds more than the sovereign spread because the rule works through the
price of capital, not through the bond.

\begin{table}[htbp]
\centering
\caption{The announced proportional rule at the impact of a risk shock}
\label{tab:proportional}
\begin{tabular}{lcccc}
\toprule
 & No backstop & Surprise purchase & Rule, delayed start & Announced rule \\
\midrule
\multicolumn{5}{l}{\emph{Default probability 1.5\% a quarter}} \\
Purchases (\% of $D$ debt) & 0 & 1.1 & 0 & 1.1 \\
Sovereign spread (bp) & 222 & 222 & 215 & 213 \\
Credit spread (bp) & 70 & 72 & 0 & 0 \\
Price of capital $Q_{D,t}$ (\%) & 0 & 0.01 & 0.34 & 0.35 \\
$D$ bank net worth $n_{D,t}$ (\%) & 0 & 0.1 & 2.5 & 2.6 \\
\midrule
\multicolumn{5}{l}{\emph{Default probability 6\% a quarter}} \\
Purchases (\% of $D$ debt) & 0 & 4.7 & 0 & 4.7 \\
Sovereign spread (bp) & 573 & 564 & 541 & 530 \\
Credit spread (bp) & 426 & 400 & 169 & 121 \\
Price of capital $Q_{D,t}$ (\%) & 0 & 0.08 & 0.66 & 0.73 \\
$D$ bank net worth $n_{D,t}$ (\%) & 0 & 0.9 & 6.3 & 7.2 \\
\bottomrule
\end{tabular}
\par\medskip
\begin{minipage}{0.92\linewidth}\footnotesize
\emph{Notes:} impact quarter, every column at the same state: the rest point of the economy
without a backstop, with the default probability raised to the value shown. The price of
capital and net worth are in percent deviation from the no-backstop column. ``Surprise
purchase'': the rule's purchase made unexpectedly, under the decision rules of the economy
without a backstop. ``Rule, delayed start'': the rule is known but buys nothing in the impact
quarter. Both are controlled counterfactuals, not policy equilibria. Spreads in bp a year;
the credit spread is $r^{wc}_{D,t}-r_{D,t}$. From its own rest point the announced rule buys
1.1\% and 5.0\% of the debt and leaves credit spreads of 0 and 169\,bp.
\end{minipage}
\end{table}

\paragraph{Expected capital prices.}
The $D$ bank's first-order condition for capital prices installed capital as
\begin{equation}
Q_{D,t}=\frac{\E_t\bigl[\Omega_{D,t,t+1}\bigl(z_{D,t+1}+(1-\delta)\,Q_{D,t+1}\bigr)\bigr]}
{\E_t\bigl[\Omega_{D,t,t+1}\bigr]R_{D,t}+\lambda_D\,\mu_{D,t}},
\label{eq:capital-price}
\end{equation}
where $z_{D,t+1}$ is the marginal product of capital, $\delta$ the depreciation rate,
$\Omega_{D,t,t+1}$ the banker's discount factor weighted by the value of a unit of net worth,
$R_{D,t}$ the gross deposit rate and $\lambda_D\mu_{D,t}$ the constraint premium. At the
1.5\% shock the announced rule raises $Q_{D,t}$ by 0.35\%. As an accounting decomposition of
\cref{eq:capital-price}, a higher expected resale value $\E_t[Q_{D,t+1}]$ contributes 0.24
points, the smaller constraint premium 0.17 points and the higher deposit rate $-0.07$ points;
expected dividends contribute nothing. At the 6\% shock the constraint premium dominates
(0.76 points, with 0.25 from the resale value and $-0.29$ from the deposit rate).

The accounting identity does not say what starts the rise. To find out, we replace one group
of next-quarter outcomes at a time with its value under the rule and re-clear the impact
quarter exactly, which is a controlled counterfactual. The largest single source is next
quarter's capital accumulation: banks expect the rule to keep the constraint looser
next quarter, so they expect more investment then, and with adjustment-cost curvature
$\xi=0.5$ the price of capital is about twenty times as elastic as the capital stock. That
expectation accounts for 0.14 of the 0.35 points at the 1.5\% shock. A lower expected path
for the $D$ deposit rate, through the terms of trade, adds 0.12 points and households'
valuation of next quarter 0.05 points. Everything in the bond market together adds only 0.02
points: the bond's continuation price, the banks' franchise value and the impact-quarter
purchase. At the 6\% shock the bond market adds 0.19 of 0.73 points. The rule therefore
raises today's price of capital mainly because it raises tomorrow's.

\paragraph{Revaluation and the constraint.}
Banks carry the undepreciated capital stock $(1-\delta)K_{D,t}$ into the quarter and value it
at the post-investment price $Q_{D,t}$, so a change in that price passes into net worth one
for one through
\begin{equation}
\Delta n_{D,t}=(1-f_D)(1-\delta)\,K_{D,t}\,\Delta Q_{D,t},
\label{eq:revaluation}
\end{equation}
where $f_D$ is the share of bankers who exit each quarter. Capital held by the banks is 6.4
times their net worth at the headline shock, so the 0.35\% price rise accounts for 2.2 of the
2.6\% rise in net worth; at the 6\% shock it accounts for 5.6 of 7.2\%. This decomposition is
an exact accounting identity. New capital is bought at $Q_{D,t}$ with deposits and adds
nothing to net worth, so nothing is counted twice. Because the multiplier is
$\mu_{D,t}=\max\{1-\E_t[\Omega_{D,t,t+1}]R_{D,t}\,n_{D,t}/\mathcal A_{D,t},\,0\}$, where
$\mathcal A_{D,t}$ is the divertible value of the bank's assets, a 1\% rise in $Q_{D,t}$
lowers it by about 0.056 before it reaches zero. Net worth rises 6.5 times as fast as
the price, while the divertible base, which is 82\% capital, rises only 0.8 times as fast. At
the headline shock the 0.35\% price rise lowers the multiplier by about 0.019, more than its
no-backstop value of 0.015, and the constraint goes slack.

The feedback through net worth would raise the price of capital by 0.65\% for every 1\% it
had already risen, but equilibrium amplifies the initial rise only by a factor of 1.3--1.4.
At the no-backstop allocation the promise raises the price banks will pay for capital by
0.27\%, against 0.35\% in equilibrium; at the 6\% shock by 0.53\%, against 0.73\%. Two
forces stop the loop. The extra investment must be financed with $D$ deposits, so the $D$
deposit rate is 26\,bp higher at the 1.5\% shock and 115\,bp higher at 6\%. A larger capital
stock tomorrow also lowers tomorrow's investment rate and with it the resale price.

\paragraph{Purchases and reserves.}
The bond purchase itself contributes little. Made as a surprise, the same purchase changes the
credit spread by about $+3$\,bp at the 1.5\% shock and $-25$\,bp at 6\%, against cuts of 70 and
305\,bp under the announced rule. A rule that is known but buys nothing in the impact quarter
delivers 96\% and 91\% of the rise in the price of capital and, at the headline shock, the
whole fall in the credit spread. At 6\% the impact-quarter purchase adds 48\,bp to the cut.
The purchase is small relative to what it would have to move. Sovereign bonds are less than a
fifth of the divertible base, and banks with a slack or barely binding constraint value a bond
close to its expected payoff, so a purchase of 1--5\% of the debt barely moves its price.

The purchase pays with reserves, which banks cannot divert. That makes the purchase loosen the
constraint as well as swapping risk. Reserves are 1.2\% and 5.0\% of the divertible base at
the two shocks. As a controlled counterfactual, making them as divertible as bonds leaves the
constraint slack at the 1.5\% shock and raises the credit spread from 169 to 194\,bp at 6\%.
Non-divertibility matters for the size of the cut at crisis levels of risk but not for its
existence.

\paragraph{Accuracy.}
Every response above clears the impact quarter exactly and takes next quarter's values from
the fitted decision rules. When the impact and following quarters are cleared jointly, the
rule's effect on the price of capital falls from 0.31\% to 0.26\% at the 1.5\% shock and from
0.63\% to 0.57\% at 6\%. The cut in the credit spread falls from 70 to 61\,bp and from 257 to
226\,bp, both measured from each economy's own rest point. The ranking of sources is
unchanged. What does not survive is the level of the spread: the fitted rules overstate next
quarter's value of a unit of bank net worth by 1--3\%, so joint clearing raises credit spreads
with and without the rule. At the headline shock the constraint then binds narrowly under the
rule ($\mu_{D,t}=0.002$, a credit spread of 10\,bp), and at 6\% the spreads rise from 426 to
516\,bp without the rule and from 169 to 290\,bp with it. We regard the size of the cut and its
source as established to within about 15\%, but not whether the constraint is exactly slack at
the headline shock, nor the level of spreads in the crisis region.
\placeholder{all numbers conditional on the Bohn coefficient of 0.25 (at 0.06 the constraint
binds only from about 5\% a quarter, so the headline shock would not engage it) and on
union-wide deposit clearing in the dynamics, which supplies the terms-of-trade term in the
deposit rate; recompute once both are decided}
